\section{Separated stress kernels and one phase average}
\label{sec:separated-short}

All phase norms use normalized measure on $\mathbb T_P$.
Write $B=\partial_\tau-y\cdot\nabla/2$,
$D_\mu=B-\nu\Delta+\mu-1/2$, and
$S_\mu=D_\mu^{-1}\mathbb P\operatorname{div}$.
The exponent $\mu$ ranges over a fixed compact subset $K$ of
$\{\operatorname{Re}\mu>-1/4\}$. Let $T$ be supported in the
radial disk $|z'|\le r_o$, and restrict receivers to $|y'|\ge r_i>r_o$.
Put $D=r_i-r_o$. The support condition is imposed on the tensor before
its divergence. We begin with smooth tensors and then complete in the
norms appearing below.

Let $G_t$ be the three-dimensional heat kernel. The heat--Leray kernel is
\[
 (K_t)_{ij}=\delta_{ij}G_t+
   \partial_i\partial_j\int_t^\infty G_a\,da.
\]
Differentiation and the substitution $a=|x|^2r$ give, off the origin,
\begin{equation}\label{sep:kernel}
 |D^j\nabla K_t(x)|\le C_j|x|^{-4-j},\qquad
 |D^j(K_t-K_0)(x)|\le C_jt|x|^{-5-j}.
\end{equation}
The second bound also follows by integrating
$\partial_tK_t=\Delta K_t$ from zero to $t$. Thus it remains valid
as $t\downarrow0$ at every separated observation point. At $t=0$,
$K_0=D^2(4\pi|x|)^{-1}$ off the origin; its local distributional
part is retained by the operator, but vanishes at a separated receiver.

The free semigroup is
\[
 e^{h/2}\bigl[e^{\nu(e^h-1)\Delta}f(\tau-h)\bigr](e^{h/2}y).
\]
In its Laplace integral set $\lambda=e^{-h/2}$ and integrate the
tensor divergence by parts. Heat scaling contributes $\lambda^3$,
the divergence contributes $\lambda$, and the semigroup amplitude
and differential contribute $\lambda^{-1}$ and $-2d\lambda/\lambda$.
Consequently the negative response $W_\mu=-S_\mu T$ is exactly
\begin{equation}\label{sep:resolvent}
 W_\mu(y,\tau)=-2\int_0^1\lambda^{2\mu+2}
  \int\nabla K_{\nu(1-\lambda^2)}(y-\lambda z)
       :T(z,\tau+2\log\lambda)\,dz\,d\lambda.
\end{equation}
The radial gap persists under every dilation:
$|y-\lambda z|\ge |y'|-\lambda|z'|\ge D$.
No temporal derivative of $T$ occurs.

\subsection{Thin sources, repair sources, and the stream cutoff}

For $b>3/(4+j)$, integrating \eqref{sep:kernel} outside the gap gives
\[
 \|1_{|x|>D}D^j\nabla K_t\|_{L^b}
       \le C_{j,b}D^{-4-j+3/b}.
\]
In Young's inequality we use $b\ge1$. Dilation of an $L^p$ source
contributes $\lambda^{-3+3/p}$, leaving the integrable power
$\lambda^{2\operatorname{Re}\mu-1+3/p}$ in
\eqref{sep:resolvent}. Minkowski and phase-translation invariance
therefore prove, whenever $1+1/q=1/p+1/b$ and
$\inf_{\mu\in K}(2\operatorname{Re}\mu+3/p)>0$,
\begin{equation}\label{sep:young}
 \|1_{\rm young}D^jW_\mu\|_{L^2_\tau L^q_y}
 \le C_{K,j,p,q}D^{-4-j+3/b}
                   \|T\|_{L^2_\tau L^p_y}.
\end{equation}
The pointwise version has $q=\infty$, $b=p'$.
In particular, for $p=22/19$,
\begin{align}
 \|\nabla W_\mu(y,\cdot)\|_{L^2_\tau}
   &\le C_KD^{-101/22}\|T\|_{L^2_\tau L^{22/19}},\label{sep:repair-point}\\
 \|1_{\rm young}\nabla W_\mu\|_{L^2_\tau L^{22/3}}
   &\le C_KD^{-46/11}\|T\|_{L^2_\tau L^{22/19}},\label{sep:repair-global}\\
 \|1_{\rm young}D^jW_\mu\|_{L^2_\tau L^2}
   &\le C_{K,j}D^{-23/11-j}\|T\|_{L^2_\tau L^{22/19}}.
 \label{sep:repair-energy}
\end{align}
Here the kernel exponents are respectively $22/3$, $11/3$, and
$11/7$. A repair tensor $T=r_b\otimes w+w\otimes r_b$ obeys
\[
 \|T\|_{L^2_\tau L^{22/19}}
 \le2\|r_b\|_{L^\infty_\tau L^{11/4}}\|w\|_{Y_2}.
\]
These estimates include the Cartesian pressure projection.

For a thin source with
$\int\|T(z,\cdot)\|_{L^2_\tau}\,dz\le R^2N_0$ and
$D\ge c_0R\sqrt\nu$, the same argument with $p=1$ gives, for
fixed $q>1$ and fixed $j$,
\begin{equation}\label{sep:thin-global}
 \|D^jW_\mu\|_{L^2_\tau L^q(\rm young)}
       \le P_\nu R^{-2-j+3/q}N_0.
\end{equation}
The polynomial $P_\nu$ depends on the fixed derivative and exponent
budgets and on the fixed fractional radial gap.

For axisymmetric no-swirl divergence-free $W$, choose a smooth radial
cutoff $\chi$ whose transition lies outside the source disk, on an
annulus of radius comparable to $D$. Its divergence-free stream cutoff is
\[
 (\mathcal C_\chi W)_r=\chi W_r,\qquad
 (\mathcal C_\chi W)_z=\chi W_z-
  \frac{\chi'(\rho)}\rho\int_\rho^\infty\rho'W_z(\rho',z)\,d\rho'.
\]
The stream primitive converges by the separated kernel bound.
The cutoff correction has a translation-invariant majorant of the form
\[
 C D^{-2}(D^2+|z-z_{\rm src}|^2)^{-1}
       1_{|y'-z'_{\rm src}|\le CD}.
\]
Indeed the radial integral of the order-four kernel is bounded by
$C(D^2+|z-z_{\rm src}|^2)^{-1}$, while $\chi'/\rho=O(D^{-2})$.
After $j$ derivatives the same estimate costs $D^{-j}$.
Its $L^b$ norm is $CD^{-4-j+3/b}$ for $b>1/2$, as follows by
scaling the two horizontal variables and the vertical integral.
Thus \eqref{sep:young}--\eqref{sep:thin-global} remain valid after
this cutoff in all exponents used here. The estimate is an inequality
for the full tensor convolution; it does not discard a cutoff source.

Put
\[
 X_0=W_\mu-W_{\mu+1},\quad X=\mathcal C_\chi X_0,
 \quad \mathcal Q=[B,\mathcal C_\chi]-\nu[\Delta,\mathcal C_\chi],
 \quad Y=\mathcal C_\chi W_\mu+\mathcal QX_0.
\]
The difference in $X_0$ inserts $1-\lambda^2$ in
\eqref{sep:resolvent}. It costs no inverse. Derivatives of $\chi$
cost $D^{-1}$, $B\chi$ is bounded, and each extra Laplacian costs
$\nu D^{-2}$. For $D\ge c_0R\sqrt\nu$ and $R\ge1$ this last
factor is bounded. Therefore the displayed estimates also hold for
$X,Y$, with finitely many increased derivative budgets.

\subsection{One phase primitive gives an extra curvature factor}

Set $\epsilon=R^{-1}$, $P=2\epsilon c$, and $x=y/R$.
Write the thin physical stress as
\[
 T(Rx,\tau)=F(x',Y,q_\epsilon(|x'|,\tau)),\qquad
 Y=x_3/\epsilon,\quad
 q_\epsilon(r,\tau)=\epsilon^{-1}\log r+\tau/(2\epsilon),
\]
where $F$ has $q$-period $c$. Define
\[
 N_j=\int(1+|Y|)^j
       \|F(x',Y,\cdot)\|_{L^2(dq/c)}\,dx'\,dY,
       \qquad j=0,1.
\]
Then its physical phase-averaged mass is $R^2N_0$.
Changing variables in \eqref{sep:resolvent} gives
\begin{equation}\label{sep:outer}
 W_\mu=-2\epsilon^2\int_0^1\lambda^{2\mu+2}
  \int\nabla K_{\nu\epsilon^2(1-\lambda^2)}
       (x-\lambda(x',\epsilon Y))
 F(x',Y,q_\epsilon(r',\tau+2\log\lambda))\,dx'\,dY\,d\lambda.
\end{equation}
Let $\bar F$ be its normalized phase mean, and set
\[
 V_{\mu,0}(x)=-2\int_0^1\lambda^{2\mu+2}
  \int\nabla K_0(x-\lambda(x',0))\bar F(x',Y)
                                  \,dx'\,dY\,d\lambda.
\]
Taylor expansion in $\epsilon Y$ and \eqref{sep:kernel} control
the mean collapse and heat errors by first transverse moments. The
strict radial gap remains valid along the entire Taylor segment.

For the oscillatory part let $J=\partial_q^{-1}$ on zero-mean
periodic functions. Fourier series gives
$\|J\|_{L^2(dq/c)\to L^2(dq/c)}\le c/(2\pi)$. Exactly in
distributions in $\lambda$,
\[
 F_\perp(q_\epsilon(r',\tau+2\log\lambda))
 =\epsilon\lambda\partial_\lambda
   JF_\perp(q_\epsilon(r',\tau+2\log\lambda)).
\]
Integrate once by parts in \eqref{sep:outer}. The boundary at
$\lambda=1$ is retained. The boundary at zero vanishes in normalized
$L^2_\tau$, since $2\operatorname{Re}\mu+3>0$ and translations of
$JF_\perp$ preserve its norm. The derivative of the kernel is
\[
 -(x',\epsilon Y)\cdot\nabla(\nabla K_t)
       -2\nu\epsilon^2\lambda\Delta(\nabla K_t),
       \qquad t=\nu\epsilon^2(1-\lambda^2).
\]
The factor $x'$ is bounded by the source radial endpoint and the
factor $Y$ by $N_1$. Every fixed spatial observation derivative is
estimated by \eqref{sep:kernel}. This proves
\begin{equation}\label{sep:mean-comparison}
 \|\nabla_yW_\mu-\epsilon^3\nabla_xV_{\mu,0}\|_{L^2_\tau}
       \le P_\nu\epsilon^4N_1
\end{equation}
on compact separated observation sets, with the stated polynomial
dependence on radial endpoints, their gap, and $c$.
No age or phase derivative of $F$ has been used.

For the full young band $s\le1/\nu$, old support $t\ge1.2/\nu$,
and $r=s^{-1/2}$ in outer coordinates, the same calculation gives
the more explicit bounds
\begin{align}
 \|D_yW_{\mu,\perp}\|_{L^2_\tau}
     &\le P_\nu\epsilon^4s^{5/2}N_1,\label{sep:phase-age}\\
 \|D_yW_\mu-\epsilon^3D_xV_{\mu,0}\|_{L^2_\tau}
     &\le P_\nu\epsilon^4(s^{5/2}+s^3)N_1.
 \label{sep:mean-age}
\end{align}
The direct differentiated kernel contributes $r^{-5}$; a transverse
collapse costs one further derivative, $r^{-6}$. In the phase
calculation the extra factor $x'$ is at most a fixed multiple of $r$;
the $Y$ term has the harmless factor $\epsilon/r$, and the heat term
has $\nu\epsilon^2/r^2\le\epsilon^2$.
These estimates hold with any fixed number of spatial observation
derivatives, with the corresponding powers of $r^{-1}$.
The stream cutoff estimates proved above extend them to $X,Y$.
Moment-norm approximation proves the statements for all $N_1$ inputs,
without assuming smoothness of the source.

\subsection{Phase suprema use a separate ordinary bound}

The extra $\epsilon$ in \eqref{sep:phase-age} is not asserted after
arbitrary phase or conormal derivatives. Instead use the exact free
equations
\[
 BX_0=\nu\Delta X_0-(\mu+1/2)X_0+W_\mu,
 \qquad
 BW_\mu=\nu\Delta W_\mu-(\mu-1/2)W_\mu
                                  -\mathbb P\operatorname{div}T.
\]
At a separated receiver the last term has an order-four kernel, and
its gradient has order five. It is bounded by precisely the ordinary
kernel estimates already proved. It does not require differentiating
$T$, but need not retain the extra phase-average factor.
For the cutoff terms use
\[
 BT_{\rm rad}W=T_{\rm rad}BW-T_{\rm rad}W,
 \qquad [B,\Delta]=\Delta,
\]
and
\[
 [B,\mathcal Q]=[B,[B,\mathcal C_\chi]]
 -\nu[\Delta,\mathcal C_\chi]
 -\nu[\Delta,[B,\mathcal C_\chi]].
\]
Thus $BX,BY$ have the same ordinary separated bounds as $X,Y$,
with a fixed increased derivative budget and the factors
$\nu D^{-2}$ already controlled. Also
$B\partial_j=\partial_jB+\partial_j/2$.

For a global cutoff field $Z$ among these vectors or tensors and
$1<q<\infty$, spatial
integration by parts gives
\[
 \left|\frac{d}{d\tau}\|Z(\tau)\|_q\right|
 \le\|BZ(\tau)\|_q+\frac3{2q}\|Z(\tau)\|_q.
\]
Choose a phase where the norm is no greater than its normalized $L^2$ average,
then integrate over an interval of length at most $P\le1$. This yields
\begin{equation}\label{sep:phase-sup}
 \|Z\|_{L^\infty_\tau L^q}
 \le(1+3P/(2q))\|Z\|_{L^2_\tau L^q}
                         +P\|BZ\|_{L^2_\tau L^q}.
\end{equation}
In particular \eqref{sep:repair-global} for $\nabla X,\nabla Y$
holds also in phase supremum, with $D\ge c_0R\sqrt\nu$ and a
fixed viscosity polynomial. The thin bounds have the same upgrade.
For weighted spatial suprema, apply the transport identity along
$y(\tau)=e^{-\tau/2}y(0)$ to
$\rho^{4+j}D^jX$ and $\rho^{4+j}D^jY$; the weight has
$B\rho^{4+j}=-(4+j)\rho^{4+j}/2$. The same argument bounds the supremum in time by the $L^2$ average. Its input is the $L^2_\tau L^\infty_y$ bound:
take the weighted spatial supremum of the separated kernel first,
then apply Minkowski in the source and dilation variables. The thin
$L^1_yL^2_\tau$ moment supplies this bound. A pointwise-in-space
$L^2_\tau$ estimate alone would not suffice.
Approximation of the source in its stated norm
passes all these estimates to rough tensors.
